%% toc: P1 3 3.1,3.2,3.6
\section{Partial orders, linear orders, and well-orders}
\label{sec:p01-w02-p01-008-orders}

The relation properties of \cref{sec:set-prerequisites-03} define strict
linear orders and strict well-orders \cite[ch.~``Sets, Relations, Functions,
and Cardinality'', sec.~``Relations and functions'', definition ``Properties
of a binary relation'']{FrozenFormalLogicOrdersChoice}. They do not define
antisymmetry, partial orders, or the reflexive (non-strict) forms of linear
orders and well-orders, and they do not say how the two forms are related.
This section supplies those definitions and proves that the strict and the
non-strict forms of each kind of order determine each other by removing or
adding the identity relation.

\emph{Inputs.} A set \(A\) and a binary relation on \(A\), that is, a subset
of \(A\times A\) (\cref{sec:set-prerequisites-03}); the clauses reflexive,
irreflexive, symmetric and transitive defined there; and the definitions of
strict linear order and strict well-order given there, which are used
unchanged. \emph{Conclusions.} The definitions of
\cref{def:p01-w02-p01-008-order-kinds}; the correspondence theorem
(\cref{thm:p01-w02-p01-008-strict-nonstrict}), which pairs strict partial
orders with partial orders, strict linear orders with non-strict linear
orders, and strict well-orders with non-strict well-orders; and the
comparison of the two readings of ``least member''.

\subsection{Definitions}
\label{sec:p01-w02-p01-008-order-defs}

\begin{hbdefinition}{def:p01-w02-p01-008-order-kinds}{Partial orders and the non-strict forms of linear orders and well-orders}
Let \(R\) be a binary relation on a set \(A\), and let
\(\mathrm{id}_A\defeq\{(a,a):a\in A\}\) be the \emph{identity relation} on
\(A\), formed by separation from \(A\times A\).
\begin{enumerate}
  \item \(R\) is \emph{antisymmetric} iff \(aRb\) and \(bRa\) imply \(a=b\)
        for all \(a,b\in A\); it is \emph{asymmetric} iff \(aRb\) implies
        that \(bRa\) fails, for all \(a,b\in A\).
  \item \(R\) is a \emph{partial order} on \(A\) iff it is reflexive,
        antisymmetric and transitive. It is a \emph{strict partial order}
        on \(A\) iff it is irreflexive and transitive.
  \item \(R\) is a \emph{non-strict linear order} on \(A\) iff it is a
        partial order on \(A\) and, for all \(a,b\in A\), \(aRb\) or
        \(bRa\).
  \item Let \(B\subseteq A\) and \(b\in B\). The member \(b\) is
        \emph{\(R\)-minimal} in \(B\) iff no \(a\in B\) satisfies \(aRb\);
        this is the condition that \cref{sec:set-prerequisites-03} calls
        being an \(R\)-least member when \(R\) is a strict linear order. The
        member \(b\) is \emph{\(R\)-least} in \(B\) iff \(bRa\) for every
        \(a\in B\).
  \item \(R\) is a \emph{non-strict well-order} of \(A\) iff it is a
        non-strict linear order on \(A\) and every nonempty
        \(B\subseteq A\) has an \(R\)-least member.
\end{enumerate}
\end{hbdefinition}

The strict linear orders and strict well-orders of
\cref{sec:set-prerequisites-03} are unchanged: a strict linear order is
irreflexive and transitive and relates exactly one way each two distinct
members, and a strict well-order is a strict linear order in which every
nonempty subset has a minimal member in the sense of item~4. ``Well-ordered''
without qualification means strictly well-ordered, as in
\cref{sec:set-prerequisites-03}. A partial order need not compare every two
members; a non-strict linear order is also called a \emph{total} order.

\subsection{The correspondence between strict and non-strict orders}
\label{sec:p01-w02-p01-008-correspondence}

\begin{hbtheorem}{thm:p01-w02-p01-008-strict-nonstrict}{proved=annex:p01-w02-p01-008-strict-nonstrict}{Strict and non-strict orders correspond}
Let \(A\) be a set.
\begin{enumerate}
  \item Every strict partial order \(R\) on \(A\) is asymmetric, and
        \(R\cup\mathrm{id}_A\) is a partial order on \(A\) with
        \((R\cup\mathrm{id}_A)\setminus\mathrm{id}_A=R\).
  \item For every partial order \(Q\) on \(A\), \(Q\setminus\mathrm{id}_A\)
        is a strict partial order on \(A\) with
        \((Q\setminus\mathrm{id}_A)\cup\mathrm{id}_A=Q\).
\end{enumerate}
Hence \(R\mapsto R\cup\mathrm{id}_A\) and \(Q\mapsto Q\setminus\mathrm{id}_A\)
are mutually inverse bijections between the set of strict partial orders on
\(A\) and the set of partial orders on \(A\). For the rest of the theorem let
\(R\) be a strict partial order on \(A\) and \(Q=R\cup\mathrm{id}_A\).
\begin{enumerate}
  \setcounter{enumi}{2}
  \item \(R\) is a strict linear order on \(A\) iff \(Q\) is a non-strict
        linear order on \(A\).
  \item If \(R\) is a strict linear order on \(A\), \(B\subseteq A\) and
        \(b\in B\), then \(b\) is \(R\)-minimal in \(B\) iff \(b\) is
        \(Q\)-least in \(B\).
  \item \(R\) strictly well-orders \(A\) iff \(Q\) is a non-strict
        well-order of \(A\).
\end{enumerate}
\end{hbtheorem}

\begin{hbmethod}{def:p01-w02-p01-008-convert}{Changing the form of an order}
Given a relation \(R\subseteq A\times A\):
\begin{enumerate}
  \item To obtain the non-strict form, first check that \(R\) is irreflexive
        and transitive. Then form \(R\cup\mathrm{id}_A\). Linearity and the
        well-order property transfer by items~3 and~5 of
        \cref{thm:p01-w02-p01-008-strict-nonstrict}.
  \item To obtain the strict form of a relation \(Q\), first check that
        \(Q\) is reflexive, antisymmetric and transitive. Then form
        \(Q\setminus\mathrm{id}_A\).
  \item A least or minimal member found in one form is, for linear orders,
        the same member in the other form (item~4).
\end{enumerate}
Side conditions: the identity relation is taken on the same set \(A\) on
which the order is declared; the checks in steps~1 and~2 are not optional,
because without them the converted relation can fail to be an order of the
other kind (\cref{def:p01-w02-p01-008-order-limits}).
\end{hbmethod}

\begin{hbexample}{def:p01-w02-p01-008-order-example}{Divisibility and the order of a finite ordinal}
\emph{A partial order that is not linear.} Let \(A=\{1,2,3,6\}\) and let
\(Q\) relate \(a\) to \(b\) iff \(a\) divides \(b\). Then \(Q\) is
reflexive, transitive and antisymmetric on \(A\) (two positive integers that
divide each other are equal), so \(Q\) is a partial order. It is not linear:
neither of \(2\) and \(3\) divides the other. Its strict form
\(R=Q\setminus\mathrm{id}_A\) is proper divisibility,
\[
  R=\{(1,2),(1,3),(1,6),(2,6),(3,6)\},
\]
which is irreflexive and transitive. In \(B=\{2,3\}\) both members are
\(R\)-minimal, and neither is \(Q\)-least; so item~4 of
\cref{thm:p01-w02-p01-008-strict-nonstrict} needs linearity.

\emph{A strict well-order and its non-strict form.} Let
\(3=\{0,1,2\}\) (\cref{thm:p01-w02-p01-008-naturals}) with
\(R=\{(0,1),(0,2),(1,2)\}\), the relation \(\in\) on \(3\). This is a strict
well-order, and
\(R\cup\mathrm{id}_3=\{(0,0),(1,1),(2,2),(0,1),(0,2),(1,2)\}\) is a
non-strict well-order: every nonempty subset has a member related to all its
members, for instance \(1\) in \(\{1,2\}\).
\end{hbexample}

\begin{hbwarning}{def:p01-w02-p01-008-order-limits}{Where the correspondence stops}
\begin{enumerate}
  \item Antisymmetry is needed for item~2 of
        \cref{thm:p01-w02-p01-008-strict-nonstrict}. On \(A=\{0,1\}\) the
        relation \(A\times A\) is reflexive and transitive but not
        antisymmetric. Removing the identity leaves
        \(\{(0,1),(1,0)\}\), which is not transitive, since it relates
        \(0\) to \(1\) and \(1\) to \(0\) but not \(0\) to \(0\).
  \item Irreflexivity is needed for item~1. The relation \(\mathrm{id}_A\)
        is transitive but, for nonempty \(A\), not irreflexive, and
        \((\mathrm{id}_A\cup\mathrm{id}_A)\setminus\mathrm{id}_A=\varnothing\)
        does not return it.
  \item Without linearity, minimal and least members differ
        (\cref{def:p01-w02-p01-008-order-example}). A strict partial order in
        which every nonempty subset has a minimal member need not be
        linear: the empty relation on \(\{0,1\}\) is an example.
  \item When \(A=\varnothing\), the empty relation is at once a strict and a
        non-strict well-order of \(A\), since \(\mathrm{id}_A=\varnothing\);
        whether an order is strict is therefore a property of a declared
        relation on a declared set, not of the relation alone.
\end{enumerate}
\end{hbwarning}

\subsection{Relationships and sources}
\label{sec:p01-w02-p01-008-orders-sources}

The definitions extend the property list of \cref{sec:set-prerequisites-03},
whose strict well-orders are those used in \cref{def:ordinal-cardinal}. The
partially ordered sets to which Zorn's lemma applies in the set-theory part
(\cref{thm:p06-zorn}) carry partial orders in the sense of
\cref{def:p01-w02-p01-008-order-kinds}, item~2.
The non-strict reading of an ordinal's order, \(\alpha\leq\beta\) iff
\(\alpha\in\beta\) or \(\alpha=\beta\), is item~1 of
\cref{thm:p01-w02-p01-008-strict-nonstrict} applied to \(\in\). The source
definition supplies the four relation properties and the two strict orders
\cite[ch.~``Sets, Relations, Functions, and Cardinality'', sec.~``Relations
and functions'', definition ``Properties of a binary
relation'']{FrozenFormalLogicOrdersChoice}. The remaining definitions, the
theorem and its proof, the procedure, the example and the limitations are
editorial.

\section{Finite ordinals, natural numbers, and countability}
\label{sec:p01-w02-p01-008-naturals}

The definitions of finite set, countable set and countably infinite set in
\cref{sec:set-prerequisites-05} quantify over a set \(\N\) whose members
\(n\) satisfy \(n=\{0,\ldots,n-1\}\) \cite[ch.~``Sets, Relations, Functions,
and Cardinality'', sec.~``Finite, countable, and uncountable sets'',
definitions ``Finite set'' and ``Countable set'']{FrozenFormalLogicOrdersChoice}.
This section constructs that set in the metalanguage: the natural numbers
are the finite ordinals, \(\N\) is the least inductive set \(\omega\), and
the identity \(n=\{0,\ldots,n-1\}\) is proved. It then proves induction on
\(\N\), the order properties of \(\N\), and the pigeonhole facts on which
the uniqueness of a finite set's size and the infinitude of \(\N\) rest.

\emph{Inputs.} The ambient classical set theory of
\cref{ch:set-prerequisites}, of which this section uses Extensionality, the
empty set, Pairing, Union and Separation, together with the Axiom of
Infinity in the form stated below; items~6 and~7 of
\cref{thm:p01-w02-p01-008-pigeonhole} also use Replacement, through the
representation theorem of \cref{sec:p01-w02-p01-008-choice}, and the
elementary ordinal facts of \cref{thm:p01-w02-p01-011-ordinals} are used
throughout. Foundation and the Axiom of Choice are not used. \emph{Conclusions.} The set \(\omega\) and its leastness
(\cref{thm:p01-w02-p01-008-omega}); the induction principle
(\cref{thm:p01-w02-p01-008-induction}); the structure of \(\N\) as an ordinal
whose members are the finite ordinals (\cref{thm:p01-w02-p01-008-naturals});
and the pigeonhole principle with its consequences for finite and countable
sets (\cref{thm:p01-w02-p01-008-pigeonhole}).

\subsection{The least inductive set}
\label{sec:p01-w02-p01-008-omega}

\begin{hbdefinition}{def:p01-w02-p01-008-inductive}{Successor and inductive set}
For a set \(x\), the \emph{successor} of \(x\) is \(S(x)\defeq x\cup\{x\}\).
A set \(I\) is \emph{inductive} iff \(\varnothing\in I\) and \(S(x)\in I\)
for every \(x\in I\). The \emph{Axiom of Infinity} of the ambient set theory
states that some inductive set exists.
\end{hbdefinition}

The successor exists by Pairing and Union, as the union of \(x\) and
\(\{x\}\). For an ordinal \(\alpha\), \(S(\alpha)\) is the successor ordinal
\(\alpha+1\) of \cref{def:ordinal-cardinal}. The condition ``\(x\) belongs
to every inductive set'' is a formula of the ambient theory, so separation
may use it.

\begin{hbtheorem}{thm:p01-w02-p01-008-omega}{proved=annex:p01-w02-p01-008-omega}{Existence and leastness of \(\omega\)}
Let \(I_0\) be an inductive set and put
\[
  \omega\defeq\{x\in I_0:x\text{ belongs to every inductive set}\}.
\]
Then \(\omega\) is inductive and \(\omega\subseteq J\) for every inductive
set \(J\). The set \(\omega\) does not depend on the choice of \(I_0\), and it
is the only inductive set included in every inductive set.
\end{hbtheorem}

\begin{hbconvention}{def:p01-w02-p01-008-naturals}{Natural numbers}
\(\N\defeq\omega\). Its members are the \emph{natural numbers}. Write
\(0\defeq\varnothing\), \(n+1\defeq S(n)\), \(1\defeq0+1\),
\(2\defeq1+1\), and so on. For \(m,n\in\N\), \(m<n\) iff \(m\in n\), and
\(m\leq n\) iff \(m<n\) or \(m=n\); this is the order of
\cref{def:ordinal-cardinal} restricted to \(\N\). For \(n\neq0\), \(n-1\)
denotes the unique \(m\in\N\) with \(m+1=n\)
(\cref{thm:p01-w02-p01-008-naturals}, item~4).
\end{hbconvention}

\begin{hbtheorem}{thm:p01-w02-p01-008-induction}{proved=annex:p01-w02-p01-008-induction}{Induction on \(\N\)}
\begin{enumerate}
  \item If \(P\subseteq\N\), \(0\in P\), and \(n+1\in P\) for every
        \(n\in P\), then \(P=\N\).
  \item Let \(\varphi(n)\) be a formula of the ambient theory, possibly with
        set parameters. If \(\varphi(0)\) holds and, for every \(n\in\N\),
        \(\varphi(n)\) implies \(\varphi(n+1)\), then \(\varphi(n)\) holds
        for every \(n\in\N\).
\end{enumerate}
\end{hbtheorem}

\subsection{The structure of \(\N\)}
\label{sec:p01-w02-p01-008-structure}

\begin{hbtheorem}{thm:p01-w02-p01-008-naturals}{proved=annex:p01-w02-p01-008-naturals}{Order structure of the natural numbers}
\begin{enumerate}
  \item Every \(n\in\N\) is an ordinal, and \(\N\) is a transitive set.
  \item \(\N\) is an ordinal. Hence \(<\) is a strict well-order of \(\N\):
        for all \(m,n\in\N\) exactly one of \(m<n\), \(m=n\), \(n<m\) holds,
        \(<\) is transitive, and every nonempty subset of \(\N\) has a
        \(<\)-least member in the non-strict sense, a \(b\) with
        \(b\leq a\) for all members \(a\).
  \item For every \(n\in\N\),
        \[
          n=\{k\in\N:k<n\}\quad\text{and}\quad
          n+1=\{k\in\N:k\leq n\}.
        \]
        Hence \(0=\varnothing\), \(1=\{0\}\), \(2=\{0,1\}\), and for every
        \(n\neq0\) the members of \(n\) are exactly the natural numbers
        \(0,\ldots,n-1\); this is the meaning of \(n=\{0,\ldots,n-1\}\) in
        \cref{sec:set-prerequisites-05}. No natural number lies strictly
        between \(n\) and \(n+1\).
  \item \(0\neq n+1\) for every \(n\in\N\); every \(n\neq0\) in \(\N\) equals
        \(m+1\) for exactly one \(m\in\N\); and \(m+1=n+1\) implies \(m=n\).
  \item \(\omega\) is a limit ordinal: it is nonzero and is not \(S(x)\) for
        any set \(x\).
\end{enumerate}
\end{hbtheorem}

\begin{hbtheorem}{thm:p01-w02-p01-008-pigeonhole}{proved=annex:p01-w02-p01-008-pigeonhole}{Pigeonhole principle}
\begin{enumerate}
  \item For every \(n\in\N\) there is no injection \(n+1\to n\).
  \item If \(m,n\in\N\) and \(m<n\), there is no injection \(n\to m\).
        Distinct natural numbers are not bijective with each other.
  \item Every \(n\in\N\) is a cardinal in the sense of
        \cref{def:ordinal-cardinal}.
  \item There is no injection \(\N\to n\) for any \(n\in\N\). Hence \(\N\)
        is not finite, \(\N\) is a cardinal, and it is an infinite cardinal
        in the sense of \cref{def:ordinal-cardinal}.
  \item If \(A\) is finite, exactly one \(n\in\N\) admits a bijection
        \(n\to A\).
  \item Every finite set is countable, and a countable set is either finite
        or countably infinite but not both.
  \item An ordinal is a finite set iff it is a natural number. Thus the
        natural numbers are exactly the \emph{finite ordinals}.
\end{enumerate}
\end{hbtheorem}

Item~6 uses the definitions of \cref{sec:set-prerequisites-05} as stated:
countable means injectable into \(\N\), and countably infinite means
bijective with \(\N\). Its last clause needs the fact that every subset of
\(\N\) is finite or bijective with \(\N\). The annex proves that fact from
the representation theorem of \cref{sec:p01-w02-p01-008-choice}, which uses
Replacement but no choice and no recursion on \(\N\).

\subsection{Procedure, example and limitations}
\label{sec:p01-w02-p01-008-naturals-use}

\begin{hbmethod}{def:p01-w02-p01-008-induction-method}{Proof by induction}
To prove \(\varphi(n)\) for all \(n\in\N\):
\begin{enumerate}
  \item State \(\varphi\) as a formula of the ambient theory, with its
        parameters fixed in advance.
  \item Prove \(\varphi(0)\).
  \item Fix \(n\in\N\), assume \(\varphi(n)\), and prove \(\varphi(n+1)\).
  \item Conclude by \cref{thm:p01-w02-p01-008-induction}, item~2.
\end{enumerate}
Side conditions: the parameters may not vary with \(n\) inside step~3; a
statement about all members of a class rather than a set is proved by
applying the procedure to a formula, never to a collection that has not been
shown to be a set. To use a least counterexample instead, apply item~2 of
\cref{thm:p01-w02-p01-008-naturals} to the set
\(\{n\in\N:\varphi(n)\text{ fails}\}\).
\end{hbmethod}

\begin{hbexample}{def:p01-w02-p01-008-naturals-example}{The first natural numbers and one pigeonhole instance}
Unwinding the convention gives
\[
  0=\varnothing,\quad
  1=\{\varnothing\},\quad
  2=\{\varnothing,\{\varnothing\}\},\quad
  3=\{\varnothing,\{\varnothing\},\{\varnothing,\{\varnothing\}\}\}.
\]
Item~3 of \cref{thm:p01-w02-p01-008-naturals} reads \(3=\{0,1,2\}\), and
\(\{k\in\N:k<3\}=\{0,1,2\}\) because \(0,1,2\in3\) and every member of
\(3\) is one of them. By item~1 of \cref{thm:p01-w02-p01-008-pigeonhole} no
map \(3\to2\) is injective. Directly: a map \(f:3\to2\) takes one of the two
values \(0,1\) at least twice among its three arguments. The set
\(A=\{\mathrm a,\mathrm b,\mathrm c\}\) of three distinct objects is finite,
by the bijection \(0\mapsto\mathrm a\), \(1\mapsto\mathrm b\),
\(2\mapsto\mathrm c\), and by item~5 no \(n\neq3\) admits a bijection onto
it.
\end{hbexample}

\begin{hbwarning}{def:p01-w02-p01-008-naturals-limits}{What the construction does and does not provide}
\begin{enumerate}
  \item The construction starts from an inductive set \(I_0\), which the
        Axiom of Infinity supplies. Without an inductive set, the definition
        of \(\omega\) has no set to separate from, and none of the other
        axioms used in this section is invoked here to supply one.
  \item Induction applies to formulas of the ambient theory and to subsets
        of \(\N\). It does not license reasoning about an informal property
        that has not been expressed by a formula.
  \item Definitions by recursion on \(\N\), such as the sequence
        \(A_0,A_1,\ldots\) in the proof of the Cantor--Schr\"oder--Bernstein
        theorem in \cref{sec:set-prerequisites-05} and the sets \(A^n\) of
        \cref{sec:set-prerequisites-02}, are not established in this
        section; they are instances of transfinite recursion with
        \(\delta=\omega\) (\cref{thm:transfinite-recursion}), proved in
        \cref{thm:p01-w02-p01-013-recursion}.
  \item The construction fixes one set-theoretic representative of the
        natural numbers. Arithmetic facts that do not depend on that
        representative are not affected, but identities such as
        \(2=\{0,1\}\) and \(1\in2\) are facts about this representative.
\end{enumerate}
\end{hbwarning}

\subsection{Relationships and sources}
\label{sec:p01-w02-p01-008-naturals-sources}

The set \(\N\) constructed here is the set used in
\cref{sec:set-prerequisites-02} for the powers \(A^n\) and in
\cref{sec:set-prerequisites-05} for finite and countable sets, and it is the
cardinal \(\aleph_0\) of \cref{sec:set-prerequisites-06}
(\cref{thm:p01-w02-p01-008-cardinal-assignment}, item~5). The ordinal facts
used in the annex are those of \cref{thm:p01-w02-p01-011-ordinals}. The
object-language counterpart in the set-theory part is the definition of
\(\omega\) in \cref{sec:zf-zfc-01}; the construction here is the same,
carried out in the metalanguage. The definitions of finite and countable
sets and the notation \(n=\{0,\ldots,n-1\}\) follow the source
\cite[ch.~``Sets, Relations, Functions, and Cardinality'', sec.~``Finite,
countable, and uncountable sets'', definitions ``Finite set'', ``Countable
set'' and ``Cardinal comparison'']{FrozenFormalLogicOrdersChoice}. The
successor notation \(S(x)\), the axiom of infinity in inductive-set form and
the separation construction of \(\omega\) follow the source's set-theory
chapter \cite[ch.~``The Axioms of ZF and ZFC'', definitions ``Infinity'' and
``The set \(\omega\)'']{FrozenFormalLogicOrdersChoice}. The theorems and
proofs of this section are editorial.

\section{Choice, well-ordering, and cardinal assignment}
\label{sec:p01-w02-p01-008-choice}

\Cref{sec:set-prerequisites-06} states the Axiom of Choice, the
well-ordering theorem, and the corollary that every set is bijective with a
unique cardinal (\cref{cor:cardinal-assignment}), and it asserts that the
order of cardinals agrees with comparison by injections. It cites a proof
for each claim \cite[ch.~``Sets, Relations, Functions, and Cardinality'',
sec.~``Choice principles used later'', definition ``Axiom of Choice'',
theorem ``Well-ordering theorem'', corollary ``Cardinal
assignment'']{FrozenFormalLogicOrdersChoice}. This section proves all of
them. The proof of the well-ordering theorem compares members of one family
of subsets by inclusion and uses neither ordinals nor recursion. Cardinal
assignment then uses the elementary ordinal facts of
\cref{thm:p01-w02-p01-011-ordinals}, which include the comparability of
ordinals, and the representation of well-orders by ordinals; none of these
uses transfinite recursion.

\emph{Inputs.} The Axiom of Choice in the form of
\cref{sec:set-prerequisites-06}: for every set \(X\) of nonempty sets there
is a function \(c:X\to\bigcup X\) with \(c(A)\in A\) for every \(A\in X\).
Ordinals and cardinals (\cref{def:ordinal-cardinal}); comparison by
injections and bijections (\cref{sec:set-prerequisites-05}); and the ambient
theory with Separation,
Power Set, Union and Replacement. \emph{Conclusions.} The representation
theorem (\cref{thm:p01-w02-p01-008-representation}, no choice); the
well-ordering theorem (\cref{thm:p01-w02-p01-008-well-ordering}); and
cardinal assignment with existence, uniqueness, agreement with injection
comparison, totality, the cardinals of \(n\) and of \(\N\), and the
converse that cardinal assignment for all sets implies the Axiom of Choice
(\cref{thm:p01-w02-p01-008-cardinal-assignment}).

\subsection{Order isomorphisms and the representation theorem}
\label{sec:p01-w02-p01-008-representation}

\begin{hbdefinition}{def:p01-w02-p01-008-isomorphism}{Order isomorphism}
Let \(\prec\) strictly well-order a set \(W\) and \(<\) strictly well-order
a set \(V\). An \emph{order isomorphism} \(f:(W,\prec)\cong(V,<)\) is a
bijection \(f:W\to V\) such that \(v\prec w\) iff \(f(v)<f(w)\), for all
\(v,w\in W\). An ordinal \(\alpha\) is ordered by \(\in\). For \(w\in W\),
the \emph{initial segment} below \(w\) is
\(\operatorname{seg}(w)=\{v\in W:v\prec w\}\), ordered by the restriction of
\(\prec\).
\end{hbdefinition}

This is the notion of order isomorphism and initial segment used in
\cref{def:p01-w02-p01-011-readings}, item~4.

\begin{hbtheorem}{thm:p01-w02-p01-008-representation}{proved=annex:p01-w02-p01-008-representation}{Every well-order is isomorphic to exactly one ordinal}
\begin{enumerate}
  \item Let \(\prec\) strictly well-order a set \(W\). There is exactly one
        ordinal \(\alpha\) with \((W,\prec)\cong(\alpha,\in)\), and exactly
        one order isomorphism from \((W,\prec)\) onto \((\alpha,\in)\).
  \item Let \(\beta\) be an ordinal and \(X\subseteq\beta\), ordered by
        \(\in\). If \((X,\in)\cong(\alpha,\in)\) for an ordinal \(\alpha\),
        then \(\alpha\leq\beta\).
\end{enumerate}
No form of the Axiom of Choice is used.
\end{hbtheorem}

\subsection{The well-ordering theorem}
\label{sec:p01-w02-p01-008-wot}

\begin{hbtheorem}{thm:p01-w02-p01-008-well-ordering}{proved=annex:p01-w02-p01-008-well-ordering}{Well-ordering}
Assume the Axiom of Choice. Then every set can be well-ordered: for every
set \(A\) there is a binary relation on \(A\) that strictly well-orders
\(A\). The proof applies the Axiom of Choice once, to the set
\(\mathcal P(A)\setminus\{\varnothing\}\).
\end{hbtheorem}

\begin{hbmethod}{def:p01-w02-p01-008-tower-method}{How the well-order is built from a choice function}
Fix a function \(c\) choosing a member \(c(X)\in X\) of each nonempty
\(X\subseteq A\). For \(B\subseteq A\), let \(B^{+}=B\cup\{c(A\setminus B)\}\)
if \(B\neq A\), and \(A^{+}=A\).
\begin{enumerate}
  \item Call a family \(\mathcal T\subseteq\mathcal P(A)\) a \emph{tower} iff
        \(\varnothing\in\mathcal T\), \(B^{+}\in\mathcal T\) for every
        \(B\in\mathcal T\), and \(\bigcup\mathcal S\in\mathcal T\) for every
        \(\mathcal S\subseteq\mathcal T\). Let \(\mathcal T_0\) be the
        intersection of all towers; it is the least tower.
  \item The members of \(\mathcal T_0\) are linearly ordered by inclusion.
  \item For \(x\in A\), let \(T_x\) be the union of the members of
        \(\mathcal T_0\) that omit \(x\). Then \(x=c(A\setminus T_x)\).
  \item Put \(x\prec y\) iff \(x\in T_y\). This is the well-order.
\end{enumerate}
Read informally, \(\prec\) lists \(A\) in the order in which \(c\) picks
elements from what remains, \(c(A)\) first. The tower construction makes this
reading exact without an ordinal-indexed sequence of choices.
\end{hbmethod}

\subsection{Cardinal assignment}
\label{sec:p01-w02-p01-008-cardinals}

\begin{hbcorollary}{thm:p01-w02-p01-008-cardinal-assignment}{proved=annex:p01-w02-p01-008-cardinal-assignment}{Cardinal assignment}
\begin{enumerate}
  \item \emph{(Existence.)} Assume the Axiom of Choice. Every set \(A\) is
        bijective with a cardinal. Specifically, \(A\) is bijective with some
        ordinal, and the least ordinal bijective with \(A\) is a cardinal.
  \item \emph{(Uniqueness.)} Distinct cardinals are not bijective with each
        other. Hence a set is bijective with at most one cardinal, and under
        the Axiom of Choice with exactly one, written \(\abs A\).
  \item \emph{(Agreement.)} Let \(\kappa\) and \(\lambda\) be cardinals and
        let \(A\) and \(B\) be sets bijective with \(\kappa\) and \(\lambda\)
        respectively. There is an injection \(A\to B\) iff
        \(\kappa\leq\lambda\); there is a bijection \(A\to B\) iff
        \(\kappa=\lambda\); and there is an injection \(A\to B\) but no
        bijection iff \(\kappa<\lambda\). Thus, under the Axiom of Choice,
        \(\abs A\leq\abs B\), \(\abs A=\abs B\) and \(\abs A<\abs B\) in
        the sense of \cref{sec:set-prerequisites-05} hold iff the
        corresponding relations hold between the cardinals \(\abs A\) and
        \(\abs B\) in the ordinal order.
  \item \emph{(Totality.)} Any two cardinals \(\kappa,\lambda\) satisfy
        \(\kappa\leq\lambda\) or \(\lambda\leq\kappa\). Under the Axiom of
        Choice, for any sets \(A,B\) there is an injection \(A\to B\) or an
        injection \(B\to A\), and every nonempty finite list
        \(\kappa_1,\ldots,\kappa_n\) of cardinals has a greatest member
        \(\max(\kappa_1,\ldots,\kappa_n)\).
  \item \emph{(Finite cardinals and \(\aleph_0\).)} For \(n\in\N\), the
        cardinal bijective with \(n\) is \(n\), and the cardinal bijective
        with \(\N\) is \(\N=\omega\), the cardinal written \(\aleph_0\). No
        choice is needed for these.
  \item \emph{(Converse.)} Without assuming the Axiom of Choice: if every set
        is bijective with some ordinal, then the Axiom of Choice holds.
\end{enumerate}
Items~2, 3 for given bijections, the first sentence of item~4, item~5 and
item~6 use no form of the Axiom of Choice.
\end{hbcorollary}

Item~1 together with item~2 is the corollary stated in
\cref{sec:set-prerequisites-06} as \cref{cor:cardinal-assignment}; item~3 is
its asserted agreement of the cardinal order with the injection-based
comparison; item~4 is the existence of the maximum used there for the Henkin
and L\"owenheim--Skolem constructions. By item~6 and
\cref{thm:p01-w02-p01-008-well-ordering}, the statement that every set is
bijective with an ordinal is equivalent to the Axiom of Choice over the
remaining axioms of the ambient theory.

\begin{hbmethod}{def:p01-w02-p01-008-cardinal-method}{Finding the cardinal of a given set}
To determine \(\abs A\):
\begin{enumerate}
  \item Find a strict well-order of \(A\); when none is evident, the
        well-ordering theorem supplies one, but it does not describe it.
  \item Find the ordinal \(\alpha\) isomorphic to it
        (\cref{thm:p01-w02-p01-008-representation}).
  \item Find the least ordinal \(\kappa\leq\alpha\) bijective with
        \(\alpha\). Then \(\abs A=\kappa\).
\end{enumerate}
Side conditions: different well-orders of \(A\) can give different ordinals
\(\alpha\) in step~2, but step~3 gives the same \(\kappa\) for all of them by
item~2 of \cref{thm:p01-w02-p01-008-cardinal-assignment}. To compare two
sets, compare their cardinals by \(\in\) or \(=\) (item~3); no further
injection need be produced.
\end{hbmethod}

\begin{hbexample}{def:p01-w02-p01-008-choice-example}{A tower on three elements, and two orders of \(\N\) with one cardinal}
\emph{The tower construction.} Let \(A=\{1,2,3\}\) and let \(c\) choose the
greatest member of each nonempty subset. Then
\(\varnothing^{+}=\{3\}\), \(\{3\}^{+}=\{2,3\}\), \(\{2,3\}^{+}=A\) and
\(A^{+}=A\), so every tower contains
\(\mathcal T_0=\{\varnothing,\{3\},\{2,3\},A\}\). This family is itself a
tower: it is closed under \({}^{+}\), and the union of any subfamily is its
largest member, or \(\varnothing\) for the empty subfamily. Hence
\(T_3=\varnothing\), \(T_2=\{3\}\), \(T_1=\{2,3\}\), and
\(3\prec2\prec1\): the order in which \(c\) picks from what remains.

\emph{Two well-orders of \(\N\).} The relation \(<\) on \(\N\) is isomorphic
to \(\omega\) by the identity. Let \(\prec\) order \(\N\) by putting \(0\)
after every other natural number and otherwise following \(<\). This is a
strict well-order, and \(0\mapsto\omega\), \(n\mapsto n-1\) for \(n\neq0\),
is an order isomorphism onto \(\omega+1=\omega\cup\{\omega\}\). So the two
well-orders give the different ordinals \(\omega\) and \(\omega+1\) in
step~2 of \cref{def:p01-w02-p01-008-cardinal-method}. Step~3 gives
\(\omega\) both times: \(\omega+1\) is bijective with \(\omega\) by
\(\omega\mapsto0\) and \(n\mapsto n+1\), and \(\omega\) is a cardinal by
item~4 of \cref{thm:p01-w02-p01-008-pigeonhole}. Hence
\(\abs\N=\aleph_0=\omega\), and \(\omega+1\) is not a cardinal.
\end{hbexample}

\begin{hbwarning}{def:p01-w02-p01-008-choice-limits}{What the choice results do and do not establish}
\begin{enumerate}
  \item The well-ordering theorem asserts existence only. The order it
        yields depends on the choice function, and nothing here produces a
        well-order of a specific set such as \(\mathcal P(\N)\) by an
        explicit description.
  \item The Axiom of Choice cannot simply be deleted from item~1 of
        \cref{thm:p01-w02-p01-008-cardinal-assignment}: by item~6, the
        conclusion for all sets implies the axiom. Whether the remaining
        axioms prove the Axiom of Choice is a separate question, which this
        section does not address.
  \item The notation \(\abs A\) has two readings: the relations
        \(\abs A\leq\abs B\), \(\abs A=\abs B\), \(\abs A<\abs B\) defined by
        injections and bijections for all sets, and the cardinal bijective
        with \(A\). Item~3 shows that the readings agree whenever the
        cardinal exists; without the Axiom of Choice only the first reading
        is available for an arbitrary set.
  \item The cardinal of a set is the least ordinal bijective with it, not
        the ordinal of a given well-order of it; the second depends on the
        order (\cref{def:p01-w02-p01-008-choice-example}).
\end{enumerate}
\end{hbwarning}

\subsection{Relationships and sources}
\label{sec:p01-w02-p01-008-choice-sources}

The section supplies the proofs that \cref{sec:set-prerequisites-06} cites
for its well-ordering theorem and for \cref{cor:cardinal-assignment}. The
counting results that rest on cardinal assignment,
\cref{thm:infinite-cardinal-finite-coding} and \cref{cor:syntax-counting},
are proved in \cref{sec:p01-w02-p01-011-finite-coding}. The ordinal facts
used are \cref{thm:p01-w02-p01-011-ordinals,thm:p01-w02-p01-011-infinite-ordinals},
and the representation of initial segments used is
\cref{thm:p01-w02-p01-011-representation}. The Part~6 statement of the
equivalent forms of the Axiom of Choice is \cref{thm:p06-zorn}, and the
implication from the Axiom of Choice to the maximal-element principle is
proved as \cref{thm:p06-ac-to-zorn}; the Ultrafilter Lemma, a weaker choice
principle, is \cref{sec:p01-w02-p01-009-ultrafilter-lemma}. The statements
of the Axiom of Choice, the well-ordering theorem and cardinal assignment,
the notation \(\abs A\) and \(\aleph_0\), and the claim that the cardinal
order agrees with injection comparison follow the source
\cite[ch.~``Sets, Relations, Functions, and Cardinality'', sec.~``Choice
principles used later'']{FrozenFormalLogicOrdersChoice}, which cites an
external text for all proofs. The proofs, the uniqueness, totality and
converse clauses, the procedures, the examples and the limitations are
editorial.
